\Def Пусть $A, B$ — множества элементов линейного пространства $L$, тогда суммой Минковского назовем $A \oplus B = \{a + b \ |\  a \in A, \ b \in B\}$.

Далее будем рассматривать в качестве $L$ плоскость или, так
называемое, $\R^2$. Под $a$ и $b$ подразумеваются радиус вектора этих точек.\newline

\Statement Сумма Минковского двух выпуклых фигур выпукла.

\Proof Рассмотрим произвольные $A_1$, $A_2 \in A$, тогда $\forall t \in [0, 1] A_t = A_1 + (A_2 - A_1)t \in A$.\newline
Аналогично $\forall t \in [0, 1] B_t = B_1 + (B_2 - B_1)t \in B$. \newline
$A_t + B_t = (A_1 + B_1) + ((A_2 + B_2) - (A_1 + B_1))t \in A \oplus B$ и это верно для произвольного $t \in [0, 1]$.\newline

\Statement Если $S, T$ - выпуклые многоугольники, то $S \oplus T$ - выпуклый многоугольник.

\Statement Если $S = \{S_1, \dots, S_n\}$ и $T = \{T_1, \dots, T_m\}$ - выпуклые многоугольники, то \begin{enumerate}
    \item В $S \oplus T$ $\mathcal{O}(n+m)$ сторон
    \item $S \oplus T = conv(\{ S_i + T_j \ | \ i \in [1, n], j\in [1, m]\})$
\end{enumerate}

\Proof Рассмотрим вектор $d \in \R^2$. Рассмотрим семейство прямых Р, для которых $d$ является нормалью. Тогда крайней точкой многоугольника в направлении $d$ назовем такую вершину $А$, что если через нее провести прямую из Р, то весь многоугольник окажется в одной полуплоскости.

Заметим, что если рассмотреть крайние вершины $S$ и $T$ в направлении $d$, то их сумма окажется тоже крайней вершиной $S \oplus T$ в данном направлении. Действительно, рассмотрим две точки, суммой которых она является. Рассмотри проекции их радиус-векторов на $d$. Тогда при взятии точек с наибольшими проекциями на $d$ мы должны были получить крайнюю точку.

Тогда можно заметить, что если рассмотреть $d$ такой, что он ортогонален одной из сторон $S$, тогда
данная сторона будет состоять из крайних точек в направлении $d$, а значит и данная сторона будет крайней в $S \oplus T$.
Из наблюдения выше рассмотрим вектора, ортогональные каждой из $|S| + |T|$ сторон, тогда они будут крайними в $S \oplus T$, откуда следует, что в $S \oplus T$ будет не больше $|S| + |T|$ вершин. Более того, выпуклость вытекает из построения того, что стороны были крайними.

Из последнего пункта последнего утверждения можно получить \textbf{наивный алгоритм} нахождения суммы Минковского:
\begin{enumerate}
    \item Пусть $C_{ij} = A_i + B_j$
    \item Построим выпуклую оболочку на $\{C_{ij}\}_{i,j=1} ^ {|A||B|}$
\end{enumerate}
\textbf{Асимптотика:} $\mathcal{O}(n^2 \cdot \log n)$, поскольку это время работы алгоритма, строящего выпуклую оболочку на $n^2$ точках.\newline

\textbf{Оптимальный алгоритм:}
\begin{enumerate}
    \item Пусть $A_0$ и $B_0$ — самые нижние (из них самые левые) точки $A$ и $B$ соответственно, $C_0 = A_0 + B_0$.
    \item Сделаем циклический сдвиг $A$ и $B$ так, чтобы обход был против часовой стрелки и начинался с соответственно $A_0$ и $B_0$.
    \item Пусть $\text{v}^A_i = \overline{A_{i - 1}A_i}, \ \text{v}^B_i = \overline{B_{i-1}B_i}$. Заметим, что массивы $\text{v}^A$ и $\text{v}^B$ отсортированы по полярному углу.
    \item Построим $\text{v}^C = Merge(\text{v}^A, \text{v}^B)$.
    \item Последовательно откладываем от $C_0$ вектора из массива $\text{v}^C$ , то есть $C_i = C_{i-1} + v^C_i$.
\end{enumerate}
Полученный набор $C_i$ -- искомый многоугольник.

\textbf{Асимптотика:} \begin{enumerate}
    \item Построение $C_0$ и препроцессинг многоугольников выполняется за $\mathcal{O}(|A| + |B|)$.
    \item Построение  $\text{v}^A$, $\text{v}^B$ выполняется за $\mathcal{O}(|A| + |B|)$.
    \item Слияние  $\text{v}^A$, $\text{v}^B$ выполняется за $\mathcal{O}(|A| + |B|)$.
    \item Построение  $C_i$ выполняется за $\mathcal{O}(|A| + |B|)$.
\end{enumerate}
Итоговое время работы: $\mathcal{O}(|A| + |B|)$.

\pagebreak